\Askhsh{13745}{4}
\begin{ylh}{1}
    \item 3.1. Είδη και στοιχεία τριγώνων
    \item 3.2. 1ο Κριτήριο ισότητας τριγώνων
    \item 3.11. Ανισοτικές σχέσεις πλευρών και γωνιών
    \item 4.2. Τέμνουσα δύο ευθειών - Ευκλείδειο αίτημα
    \item 5.4. Ρόμβος
    \item 5.6. Εφαρμογές στα τρίγωνα.
\end{ylh}
% ===================================================================
%  Αρχείο: 13745.tex (Fragment)
%  Αυτόματη μετατροπή μέσω doc2latex.py
% ===================================================================

ΘΕΜΑ 4

Δίνεται ισοσκελές τρίγωνο ΑΒΓ, το μέσο Μ της βάσης ΒΓ και τυχαίο εσωτερικό σημείο Δ στη βάση του.

\begin{alist}
\item Αν από το μέσο Μ φέρουμε παράλληλες προς τις πλευρές ΑΒ και ΑΓ του τριγώνου, που τις τέμνουν στα σημεία Ε και Ζ αντίστοιχα να αποδείξετε ότι:

.}
\end{alist}
\begin{alist}
\item
  ΜΕ = ΜΖ. \monades{6}
\item
  Το ΑΕΜΖ είναι ρόμβος με περίμετρο ίση με 2ΑΒ. \monades{7}

\item Αν πάρουμε τυχαίο εσωτερικό σημείο Δ στο ευθύγραμμο τμήμα ΒΓ, διαφορετικό από το μέσο Μ, και φέρουμε τις παράλληλες προς τις πλευρές ΑΒ και ΑΓ του τριγώνου, που τις τέμνουν στα σημεία Κ και Λ αντίστοιχα, τότε:

.}
\end{alist}
\begin{alist}
\item
  Ποιο είναι το είδος του τετράπλευρου ΑΚΔΛ;
\item
  Να συγκρίνετε την περίμετρο του τετράπλευρου ΑΚΔΛ με την περίμετρο του ρόμβου ΑΕΜΖ του ερωτήματος α ii) και να διατυπώστε λεκτικά το συμπέρασμα που προκύπτει.

\monades{12}
\end{alist}
